% Generated draft: requires mathematical and source review.
\documentclass[11pt]{article}
\usepackage[margin=25mm]{geometry}
\usepackage{fontspec}
\setmainfont{Latin Modern Roman}
\usepackage{amsmath,amssymb,mathtools,unicode-math}
\setmathfont{Latin Modern Math}
\usepackage{enumitem,needspace,xurl,hyperref}
\setlength{\parindent}{0pt}
\setlength{\parskip}{5pt}
\setlength{\emergencystretch}{3em}
\newcommand{\N}{\mathbb{N}}
\newcommand{\Z}{\mathbb{Z}}
\newcommand{\Q}{\mathbb{Q}}
\newcommand{\R}{\mathbb{R}}
\newcommand{\C}{\mathbb{C}}
\newcommand{\F}{\mathbb{F}}
\newcommand{\A}{\mathbb{A}}
\newcommand{\PP}{\mathbb{P}}
\newcommand{\E}{\mathbb{E}}
\newcommand{\eps}{\varepsilon}
\newcommand{\dd}{\,\mathrm{d}}
\begin{document}
\section*{UnsolvedMath ID: 30006640}
\textbf{Source problem code:} catalogue entry
\subsection{Definitions and mathematical statement}
The supplied record describes a research programme rather than a fully specified conjecture. Its objects are thermodynamical probability manifolds arising from reciprocal relations in master equations, but it does not specify the master equations, reciprocal relations, metric construction, or admissible class of manifolds.

For notation in the finite-state case, let \(n\geq 2\) and write
\[
\Delta_{n-1}=\left\{p\in\R^n:p_i\geq0\text{ for all }i,\ \sum_{i=1}^n p_i=1\right\},\qquad
\Delta_{n-1}^{\circ}=\{p\in\Delta_{n-1}:p_i>0\text{ for all }i\}.
\]
These definitions do not settle whether the intended geometry is on the interior, on boundary strata, or on another probability manifold. If a smooth probability manifold \(M\) and a Riemannian metric \(g\) have been specified, let \(\nabla\) be its Levi--Civita connection. With the explicitly chosen convention
\[
R(X,Y)Z=\nabla_X\nabla_YZ-\nabla_Y\nabla_XZ-\nabla_{[X,Y]}Z,
\]
the sectional curvature of a two-dimensional tangent plane \(\sigma=\operatorname{span}\{u,v\}\subset T_pM\) is
\[
K_p(\sigma)=\frac{g_p(R(u,v)v,u)}{g_p(u,u)g_p(v,v)-g_p(u,v)^2},
\]
where \(u,v\) are linearly independent. This notation interprets the requested Riemannian curvature as the curvature tensor; the source does not specify a curvature convention.

\emph{Requested research programme.} For the intended class of thermodynamical probability manifolds arising from master equations, compute geometric quantities, including the Riemannian curvature tensor and sectional curvatures; develop the associated network analysis of curvature tensors; and investigate how these computations can be used to design algorithms for sampling distributions supported on simplices. For a target probability measure \(\mu\) on a finite simplex, sampling means producing random simplex-valued outputs with law \(\mu\), or with an explicitly controlled approximation to that law. The supplied source does not specify which target measures are admissible, what approximation criterion is required, or what algorithmic guarantees are sought.

The class of master equations, its associated network and metric, boundary behaviour, and the meaning of network analysis remain unresolved. In particular, the record does not support replacing this programme by a universal existence or convergence theorem.
\subsection{Short English statement}
Determine curvature for the intended probability geometries arising from master equations, relate it to their networks, and investigate its use in sampling probability measures on simplices. The equations, metrics, boundary conditions, target distributions, and required sampling guarantees are unspecified, so this remains a research programme requiring mathematical clarification.
\subsection{Sources}
\begin{itemize}
\item Supplied local UnsolvedMath registry record, permanent ID \(30006640\), title \emph{Geometry and Sampling on Probability Manifolds}, exact source problem code \texttt{catalogue entry}; collection: Oberwolfach Reports Open Problems, proposed year 2026.
\item \emph{Flows on Measure Spaces and Applications in Machine Learning} (2026), Oberwolfach Reports, pp.\ 795--884, DOI \texttt{10.4171/owr/2026/13}: \url{https://doi.org/10.4171/owr/2026/13}. This is also the supplied literature reference labelled ``originating Oberwolfach report.'' The original statement associates the problem with Wuchen Li and the presentation \emph{Geometric calculations on probability manifolds from reciprocal relations in Master equations}; the registry supplies no separately identified proposer.
\item Li, \emph{Geometric calculations on probability manifolds}: \url{https://arxiv.org/abs/2504.19368}.
\item Li, \emph{Transport information geometry I}: \url{https://arxiv.org/abs/1803.06360}.
\item Literature assessment dated exactly \texttt{2026-08-22}: the supplied curation labels the problem partially solved and states that Riemannian, sectional, Ricci, and scalar curvature calculations are available on several probability manifolds, while a general algorithmic theory for sampling simplex-supported distributions is not established. This is supplied dated curation derived from the Oberwolfach material, not independently verified here; its time-sensitive status and cited claims require independent verification.
\item The registry labels the statement \texttt{corrected\_verified} and says it was reconstructed and checked against the official Oberwolfach report. That verification claim is supplied provenance, not an independent verification by this draft.
\end{itemize}
\end{document}
